\documentclass[12pt]{article}

%% page layout
\textheight 9in 
\textwidth 6.5in
\topmargin -0.5in
\oddsidemargin 0in
\evensidemargin 0in
%\hoffset=-0.5in
\voffset=-0.25in

%% packages
\usepackage{scrtime} % for \thistime (this package MUST be listed first!)
\usepackage{amsmath} % essential for cases environment
\usepackage{amsthm} % for theorems and proofs
\usepackage{amsfonts} % mathbb
\usepackage{graphics,graphicx}
\usepackage{multirow} % fancy tables
\usepackage{wasysym} % circle symbols (including half-filled circles)
\usepackage{enumerate} % fancier enumeration (e.g., a,b,c, ...)
%\usepackage{xcolor}
\usepackage{color}
\newcommand{\de}[1]{{\color{red}{\bfseries DE:} #1}}

%% for solutions to multiple choice questions:
\newcommand{\correct}{{\color{blue}\fbox{\color{red}\checkmark} }}

%% macros
\newcommand{\reals}{\mathbb{R}}
\newcommand{\term}[1]{{\bfseries\slshape #1}}
\newcommand{\Ker}{{\text{Ker}\,}}
\newcommand{\Range}{{\text{Range}\,}}
\newcommand{\diag}{{\text{diag}}}
\newcommand{\alg}{{\text{alg}}}
\newcommand{\geom}{{\text{geom}}}
\newcommand{\norm}[1]{\left\|#1\right\|}
\newcommand{\abs}[1]{\left|#1\right|}
\newcommand{\R}{{\cal R}}
\newcommand{\G}{{\cal G}}
\newcommand{\eps}{\varepsilon}
\newcommand{\B}{\cal B}
\newcommand{\Tinf}{T_\textrm{inf}}
\newcommand{\Shat}{{\hat{S}}}
\newcommand{\Ihat}{{\hat{I}}}
\newcommand{\ie}{\emph{i.e., }}
\newcommand{\eg}{\emph{e.g., }}
\newcommand{\Rlogo}{\protect\includegraphics[height=2ex,keepaspectratio]{images/Rlogo.pdf}\xspace}
\newcommand{\XPPAUT}{\texttt{XPPAUT}\xspace}
\newcommand{\etal}{\textit{et al}.\xspace}
\newcommand\emphblue[1]{\emph{\color{blue}#1}}

%%%%%%%%%%%%%%%%%%%%%%%%%%%%
%% from feverpreamble.tex %%
%%%%%%%%%%%%%%%%%%%%%%%%%%%%

\usepackage{amssymb,latexsym,amsmath,setspace}
%%\usepackage[colorlinks,linkcolor=blue]{hyperref}
\usepackage[colorlinks=true,allcolors=blue]{hyperref}
%%\usepackage[colorlinks]{hyperref}
\usepackage{xspace}
\usepackage{graphics,graphicx}
\usepackage{subfigure}
\usepackage{lineno}
\usepackage{fancyhdr}
\usepackage[english]{babel}  %% for texi2dvi ~ bug
\usepackage[normalem]{ulem}
\usepackage{tikz} % http://www.texample.net/tikz/examples/tikzdevice-demo/
  % N.B. version 0.6.3 of tikzDevice from Rforge is required!!
%% improve figure caption typsetting:  (see ~/tex/caption.pdf for manual)
\usepackage[footnotesize,bf]{caption}
\usepackage{placeins} % \FloatBarrier

% citation macros
\newcommand{\citen}[1]{\cite{#1}}

% comment macros
\usepackage{color}
\newcommand{\comment}[3]{\textcolor{#1}{\textbf{[#2: }\textit{#3}\textbf{]}}}
\newcommand{\david}[1]{\comment{blue}{DE}{#1}}
\newcommand{\paul}[1]{\comment{cyan}{PA}{#1}}
\newcommand{\ben}[1]{\comment{magenta}{BB}{#1}}
\newcommand{\needref}{{\textcolor{red}{[NEED REF]}}}
\newcommand{\TBD}{{\textcolor{red}{{\bf TBD}}}}
% remove comments
%\renewcommand{\comment}[3]{\relax}
% remove section headings
%\def\subsubsection*#1{\relax\nobreak}

% other macros
\newcommand{\avg}[1]{{\left\langle#1\right\rangle}}
\newcommand{\var}[1]{\textrm{var}\left(#1\right)}
\newcommand{\sem}[1]{\textrm{sem}\left(#1\right)}
\newcommand{\natinf}{{\mathcal I}}
\newcommand{\find}{f_{\textrm{i}}}
\newcommand{\fpop}{f_{\textrm{p}}}
\newcommand{\logit}{\textrm{logit}}
\newcommand{\sign}{\textrm{sign}}
\newcommand{\logistic}{\textrm{logistic}}
\def\AJE{{\it American Journal of Epidemiology\/}}
\newcommand{\code}[1]{{\tt #1}}
\newcommand{\magcode}[1]{{\tt\color{magenta}#1}}
\newcommand{\redcode}[1]{{\tt\color{red}#1}}
\newcommand{\blackcode}[1]{{\tt\color{black}#1}}

%%%%%%%%%%%%%%%%%%%
%% JOURNAL NAMES %%
%%%%%%%%%%%%%%%%%%%
\def\PNAS{PNAS}
\def\JAMA{JAMA}
\def\BMB{{\it Bulletin of Mathematical Biology\/}}

% references
\newcommand{\eref}[1]{Equation~\eqref{E:#1}}
\newcommand{\fref}[1]{Figure~\ref{F:#1}}
\newcommand{\tref}[1]{Table~\ref{T:#1}}
\newcommand{\sref}[1]{\S\ref{S:#1}}
% other macros
\newcommand{{\Reff}}{{\mathcal{R}}_{\rm eff}}
\newcommand{{\Sinit}}{S_{\rm init}}
\newcommand{\supp}{Supplementary Information}
\newcommand{\StoppedHere}{\bigskip\bigskip{\textcolor{red}{\hrule\centerline{\bfseries STOPPED HERE}\hrule}}\bigskip\bigskip}
\newcommand{\colvec}[2]{\begin{pmatrix}#1\\#2\end{pmatrix}}
\newcommand{\diagmat}[3]{\begin{pmatrix}#1&0&0\\0&#2&0\\0&0&#3\end{pmatrix}}

\newcommand{\solution}[1]{{\hfill\break\vspace{-0.5\baselineskip}\break\color{blue}\emph{Solution: }#1}}
\newcommand{\tr}{\text{tr}}

\newcommand{\thickredline}{\bigskip{\color{red}\hrule height 5pt}\bigskip}

%% for assignment 3:
\newtheorem{theorem}{Theorem}
\newtheorem{remark}{Remark}
\newcommand{\openset}{{\mathcal O}}
\newcommand{\C}{{\mathcal C}}

% Journal Names
% -------------
% \def\MNRAS{{\it Mon.\ Not.\ R.\ astr.\ Soc.\/}}
\def\MNRAS{{\it Monthly Notices of the Royal Astronomical Society\/}}
% \def\ApJ{{\it Astrophys.~J.\/}}
\def\ApJ{{\it The Astrophysical Journal\/}}
% \def\Interface{{\it J.\ R.\ Soc.\ Lond.\/ \rm Interface}}
\def\Interface{{\it Journal of the Royal Society of London, \rm Interface}}
% \def\ProcA{{\it Proc.\ R.\ Soc.\ Lond.\/ \rm A}}
\def\ProcA{{\it Proceedings of the Royal Society of London, Series\/ \rm A}}
% \def\ProcB{{\it Proc.\ R.\ Soc.\ Lond.\/ \rm B}}
\def\ProcB{{\it Proceedings of the Royal Society of London, Series\/ \rm B}}
% \def\TransB{{\it Phil.\ Trans.\ R.\ Soc.\ Lond.\/ \rm B}}
\def\TransB{{\it Philosophical Transactions of the Royal Society of London, Series\/ \rm B}}
\def\TREE{{\it Trends in Ecology and Evolution\/}}
% \def\JTB{{\it J.\ theor.\ Biol.\/}}
\def\JTB{{\it Journal of Theoretical Biology\/}}
% \def\BE{{\it Behav.\ Ecol.\/}}
\def\BE{{\it Behavioural Ecology\/}}
% \def\BES{{\it Behav.\ Ecol.\ Sociobiol.\/}}
\def\BES{{\it Behavioural Ecology and Sociobiology\/}}
% \def\BJLS{\it Biol.\ J.\ Linn.\ Soc.\/}}
\def\BJLS{{\it Biological Journal of the Linnean Society\/}}
\def\Nature{{\it Nature\/}}
\def\Science{{\it Science\/}}
\def\Lancet{{\it The Lancet\/}}
\def\LancetID{{\it Lancet Infectious Diseases\/}}
% \def\PNAS{{\it Proc.\ Natl.\ Acad.\ Sci.\ USA}}
\def\PNAS{{\it PNAS -- Proceedings of the National Academy of Sciences of the
U.S.A.}}
\def\PLoSMed{{\it PLoS Medicine\/}}
\def\PLoSCB{{\it PLoS Computational Biology\/}}
\def\EID{{\it Emerging Infectious Diseases\/}}
\def\BMB{{\it Bulletin of Mathematical Biology\/}}
\def\AJE{{\it American Journal of Epidemiology\/}}
\def\TPB{{\it Theoretical Population Biology\/}}
\def\JRSS{{\it Journal of the Royal Statistical Society\/}}
\def\JRSSB{{\it Journal of the Royal Statistical Society Series B\/}}
\def\IRV{{\it Influenza and Other Respiratory Viruses\/}}
\def\JAMA{{\it JAMA -- Journal of the American Medical Association\/}}
\def\JGLR{{\it Journal of Great Lakes Research\/}}
\def\NEJM{{\it New England Journal of Medicine\/}}
\def\JMB{{\it Journal of Mathematical Biology\/}}
\def\Interface{{\it Journal of the Royal Society Interface\/}}
\def\THEE{{\it Theoretical Ecology\/}}
\def\Annals{{\it Annals of Internal Medicine\/}}
\def\PRL{{\it Physical Review Letters\/}}
\def\PRX{{\it Physical Review X\/}}
\def\NAMS{{\it Notices of the American Mathematical Society\/}}

%% underline with smash through:
\newcommand*{\undersmash}[1]{\underline{\smash{#1}}}

%% referring to TeX macros
\newcommand\ttbackslash{{\tt\char`\\}}
\newcommand{\macro}[1]{{\tt\ttbackslash#1}}

%%%%%%%%%%%%%%%%%%%%%%%%%%%%%%%%%%%%%%%%%%%%%%%%%%%%%%%
%% QUESTIONS FOR MATH 4MB/6MB ASSIGNMENTS.           %%
%% The question texts are used in several documents: %%
%% assignment, solutions, template,                  %%
%% hence it is better to load them from this file.   %%
%%%%%%%%%%%%%%%%%%%%%%%%%%%%%%%%%%%%%%%%%%%%%%%%%%%%%%%

%% \section{Analysis of the SI model}

\newcommand{\SIanalIntro}{%
The SI model can be written
%
\begin{equation}\label{E:SI}
  \frac{dI}{dt} = \beta I(N - I) \,,
\end{equation}
%
where $I$ denotes prevalence and $N=S+I$ is the total population size.
}

\newcommand{\SIanalQa}{%
Prove that the endemic equiliibrium (EE) is a globally asymptotically stable (GAS) equilibrium by finding an appropriate Lyapunov function.  Note that ``global'' here refers to all biologically relevant initial conditions except the (unstable) disease free equilibrium (DFE).  \\
\emph{Hint:} Lyapunov functions often look paraboloidal. \\
\emph{Note:} Notions of stability and Lyapunov functions were discussed in Math 3F03 Lecture 27 in 2013 (\url{http://www.math.mcmaster.ca/earn/3F03}).
}
\newcommand{\SIanalQb}{%
In class we proved only stability of the EE, not asymptotic stability.  Prove GAS ``directly'' in two distinct ways: 
}
\newcommand{\SIanalQbi}{%
find the exact solution of the model and take the limit as $t\to\infty$, and conclude that every solution that starts in the interval $(0,N)$ converges to the EE (this approach works only in situations where you can find the exact solution);
}
\newcommand{\SIanalQbii}{%
given $\epsilon>0$, prove that for any $I(0)\in(0,N)\ \exists t<\infty$ such that $I(t)\in[N-\epsilon,N)$ and use this to establish GAS. (Do not use your exact solution in this part; the point is to use an approach that also works for models that cannot be solved exactly.)
}

%% \section{Analysis of the basic SIR  model}

\newcommand{\basicSIRanalIntro}{%
The basic SIR model is specified by the following system of differential equations.
\begin{subequations}
\begin{align}
\frac{dS}{dt} &= -\R_0 SI \\
\noalign{\vspace{5pt}}
\frac{dI}{dt} &= \R_0 SI - I \\
\noalign{\vspace{5pt}}
\frac{dR}{dt} &= I
\end{align}
\end{subequations}
The state variables $S$, $I$ and $R$ are the proportions of the population that are susceptible, infectious and removed, respectively.  The parameter $\R_0$ is the basic reproduction number.  The time unit has been chosen to be the mean infectious period for convenience.
}

\newcommand{\basicSIRanalQa}{%
A quantity of some practical importance is the \term{peak prevalence} of disease in the population, \emph{i.e.,} the maximum proportion of the population that is simultaneously infected.  Find an exact expression for the peak prevalence, given initial conditions ($S_0,I_0$).  Why might a public health official want to know this quantity?
}

\newcommand{\basicSIRanalQb}{%
It would be helpful to have an analytical expression for the solution of the model.  Most valuable would be a formula for $I(t)$, which is most closely related to time series data.  You probably will not find a formula for $I(t)$ (extra credit if you do!!) but it is definitely possible to find an exact expression that relates $R$ (proportion removed) and $t$ (time).
}

\newcommand{\basicSIRanalQbi}{%
Find such an expression.  \emph{Hint:} Combine the equations for $dS/dt$ and $dR/dt$ into one equation that can be solved for $S$ as a function of $R$.  Then recall that $S+I+R=1$ and use the $dR/dt$ equation again.  \emph{\underline{Note}:} You will end up with an expression for $t$ as a function of $R$, not $R$ as a function $t$.
}

\newcommand{\basicSIRanalQbii}{%
Use your expression for $t(R)$ to find an expression for the time at which peak prevalence will occur.  Why might this be useful?
}

\newcommand{\basicSIRanalQbiii}{%
How could your expressions be used to compare with the time series for pneumonia and influenza in Philadelphia in 1918?  (Don't actually do it; just clearly explain your thinking including any assumptions you are making.)  Would you advise your assistant who just graduated with a degree in math and biology to do this (to help you prepare your report for the public health agency)?  Why or why not?
}

\newcommand{\basicSIRanalQbiv}{%
Is it possible to find an exact analytical expression for $t$ as a function $S$?
}

\newcommand{\basicSIRanalQc}{%
Prove that all solutions of the basic SIR model approach $I=0$ asymptotically, and explain why this makes biological sense.  \emph{Hint:} Is the function $L(S,I)=I$ a Lyapunov function?  Read the \hyperlink{NotesLyapFuns}{Notes on Lyapunov functions} below.
}

\newcommand{\basicSIRanalQd}{%
Find and classify the stability of all \undersmash{equilibria} of the basic SIR model.
}

%%\section*{Notes on Lyapunov functions}\hypertarget{NotesLyapFuns}{}

\long\def\NotesOnLyapunovFunctions{%
Consider Lyapunov's Stability Theorem as stated in \href{http://lalashan.mcmaster.ca/theobio/3F03/images/3/3b/3fl28_2013.pdf}{Math 3F03 Lecture 28 in 2013}:
%%\href{../../../3F03/lectures/3fl28_2013.pdf}{Math 3F03 Lecture 28 in 2013}:
\begin{theorem}[Lyapunov's Direct Method]\label{Th:LyapFun}
  Consider an equilibrium $X_*$ of $X'=F(X)$ and an open set $\openset$ containing $X_*$.   If $\exists$ a differentiable function $L:{\openset}\to\reals$ such that
  \begin{enumerate}[(a)]
  \item $L(X_*)=0$\quad and\quad $L(X)>0$\quad $\forall X\in {\openset}\setminus\{X_*\}$\qquad(\textcolor{blue}{$L$ positive definite on $\openset$})
 \item $\dot{L}(X)\le0$\quad $\forall X\in {\openset}\setminus\{X_*\}$\qquad (\textcolor{blue}{$\dot L$ negative semi-definite on $\openset$})
  \end{enumerate}
  then $X_*$ is stable and $L$ is called a {\bfseries Lyapunov function}.  If, in addition,
  \begin{enumerate}[(c)]
  \item $\dot{L}(X)<0$\quad $\forall X\in {\openset}\setminus\{X_*\}$ \qquad (\textcolor{blue}{$\dot L$ negative definite on $\openset$})
  \end{enumerate}
 then $X_*$ is asymptotically stable and $L$ is called a {\bfseries strict Lyapunov function}.
\end{theorem}
%%This theorem does not help us at any of the continuum of DFEs if our putative Lyapunov function is $L(S,I)=S+I$ or $L(S,I)=I$, since condition (a) fails.  However, 
Theorem~\ref{Th:LyapFun} can be generalized for analysis of stability of sets more complicated than isolated equilibria, such as periodic orbits or line segments.  If you think through the proof of the theorem above (\eg \cite[\S9.2, theorem stated on p.\,193 and proved on p.\,196]{Hirs+13}), you should be able to convince yourself that the proof still works if the equilibrium $X_*$ is replaced by any \emph{closed forward-invariant set} (often simply called a \emph{closed invariant set}).  This observation allows us to state the following more general theorem.

\begin{theorem}[Lyapunov's Direct Method for Closed Invariant Sets]\label{Th:LyapFunGen}
Consider a closed invariant set $\C$ of $X'=F(X)$ and an open set $\openset$ containing $\C$.   If $\exists$ a differentiable function $L:\openset\to\reals$ such that
  \begin{enumerate}[(a)]
  \item $L(X)=0\ \forall X\in\C$\quad and\quad $L(X)>0$\quad $\forall X\in\openset\setminus\C$\qquad(\textcolor{blue}{$L$ positive definite on $\openset$})
 \item $\dot{L}(X)\le0$\quad $\forall X\in\openset\setminus\C$\qquad (\textcolor{blue}{$\dot L$ negative semi-definite on $\openset$})
  \end{enumerate}
  then $\C$ is stable and $L$ is called a {\bfseries Lyapunov function}.  If, in addition,
  \begin{enumerate}[(c)]
  \item $\dot{L}(X)<0$\quad $\forall X\in\openset\setminus\C$\qquad (\textcolor{blue}{$\dot L$ negative definite on $\openset$})
  \end{enumerate}
then $\C$ is asymptotically stable and $L$ is called a {\bfseries strict Lyapunov function}.
\end{theorem}
Note in the above theorems that open sets are defined relative to the subset of interest; in our case this subset is $\Delta=\{(S,I):S\ge0,\, I\ge0,\, S+I\le1\}$, not all of $\reals^2$.  An open set of $\Delta$ is a set of the form $U\cap\Delta$ where $U$ is an open set of $\reals^2$.  (These sets are said to be open in the \emph{\bfseries relative topology} on $\Delta$.)  In particular, note that $\Delta$ is \emph{open} as a subset of itself, in spite of the fact that it is \emph{not open} as a subset of $\reals^2$, whereas $\Delta$ is closed in both the relative topology on $\Delta$ and the usual topology on $\reals^2$.
}

%%\bibliographystyle{vancouver}
%%\bibliography{4mba1_2018}
 % questions

%%%%%%%%%%%%%%%%%%%%%%%%%%%%%%%%%%%
%% FANCY HEADER AND FOOTER STUFF %%
%%%%%%%%%%%%%%%%%%%%%%%%%%%%%%%%%%%
\usepackage{fancyhdr,lastpage}
\pagestyle{fancy}
\fancyhf{} % clear all header and footer parameters
%%%\lhead{Student Name: \theblank{4cm}}
%%%\chead{}
%%%\rhead{Student Number: \theblank{3cm}}
%%%\lfoot{\small\bfseries\ifnum\thepage<\pageref{LastPage}{CONTINUED\\on next page}\else{LAST PAGE}\fi}
\lfoot{}
\cfoot{{\small\bfseries Page \thepage\ of \pageref{LastPage}}}
\rfoot{}
\renewcommand\headrulewidth{0pt} % Removes funny header line
%%%%%%%%%%%%%%%%%%%%%%%%%%%%%%%%%%%

\begin{document}

\begin{center}
{\bfseries Mathematics 4MB3/6MB3 Mathematical Biology\\
\smallskip
2018 ASSIGNMENT {\color{blue}1}}\\
\medskip
\underline{\emph{Group Name}}: \texttt{{\color{blue}The Infective Collective}}\\
\medskip
\underline{\emph{Group Members}}: {\color{blue} Aurora Basinski-Ferris, Michael Chong, Sang Woo Park, Daniel Presta}
\end{center}

\section{Analysis of the SI model}

\SIanalIntro
\begin{enumerate}[(a)]
\item \SIanalQa
  
  {\color{blue}
    \begin{proof}
      In order to prove that the endemic equilibrium, given by $I_\ast = N$, is globally asymptotically stable, we must first find an appropriate Lyapunov function. First, let 
\begin{equation*}
\Delta = \{I \,:\, 0 \leq I \leq N\} \subset \mathbb{R}
\end{equation*}
represent an open set consisting of the biologically relevant region, containing $I_\ast = N$. Then, consider the the following $C^1$ function $L: \Delta \to \mathbb{R}$ given by
\begin{equation*}
L(I) = (I^2 - N^2)^2.
\end{equation*}
Since $L(N) = 0$ and $L(I) > 0$ for all $I \in \Delta \setminus \{I_\ast\}$, we say that $L$ is positive definite on $\Delta$. Additionally, observe that
\begin{equation*}
\begin{aligned}
\dot{L}(I) &= \frac{dL}{dI} \frac{dI}{dt} \\
&= 2(I^2 - N^2)(2 I)(\beta I)(N-I) \\
&= 4 \beta I^2 (I^2 - N^2)(N-I)\\
& = - 4 \beta I^2  (I+N) (N-I)^2
\end{aligned}
\end{equation*}
is negative for all $I \in \Delta \setminus \{I_\ast\}$. As a result, $L$ is negative definite on $\Delta$. By Lyapunov's Direct Method, $L$ is a strict Lyapunov function, and $I_\ast=N$ is asymptotically stable.

In order to prove global asymptotic stability, we observe that as the magnitude of $I$ gets arbitrarily large, $L(I)$ also gets arbitrarily large. In other words, $L(I) \to \infty$ as $|I| \to \infty$, and $L$ is thus radially unbounded. Since $L(I)$ is radially unbounded, and $L(I) < 0$ $\forall \ I \in \Delta \setminus \{I_\ast\}$, it follows by LaSalle's Invariance Principle that the endemic equilibrium is globally asymptotically stable.
    \end{proof}
  }
  
\item \SIanalQb
  \begin{enumerate}[(i)]
  \item \SIanalQbi
    
    {\color{blue}
      \begin{proof}
        In order to find an exact solution of the model, we must first solve the separable ordinary differential equation given by
\begin{equation*}
\frac{dI}{I(N-I)} = \beta dt.
\end{equation*}
Integrating both sides, we obtain the following:
\begin{equation*}
\begin{aligned}
\int \frac{dI}{I(N-I)} &= \int \beta dt \\
- \frac{1}{N} \int \left(\frac{1}{I-N} - \frac{1}{I}\right) dI &= \int \beta dt\\
\frac{1}{N} \ln{|I|} - \frac{1}{N} \ln{|I-N|} + C_2 &= \beta t + C_1.
\end{aligned}
\end{equation*}
Let $C_3 = C_1 - C_2$. Then, we obtain the following equation:
\begin{equation*}
\ln{\left(\frac{I}{I-N}\right)} = N(\beta t + C_3).
\end{equation*}
To further simplify, let $C = e^{NC_3}$ to obtain the following solution:
\begin{equation}\label{genI}
\begin{aligned}
I &= \frac{NCe^{N \beta t}}{Ce^{N \beta t} - 1}.
\end{aligned}
\end{equation}

Assuming an initial condition of $I(0) = I_0$, we can then solve for $C$:
\begin{equation*}
\begin{aligned}
I_0 &= \frac{NCe^{N\beta(0)}}{Ce^{N\beta(0)-1}} \\
I_0 &= \frac{NC}{C-1} \\
C I_0 - NC &= I_0 \\
C &= \frac{I_0}{I_0 - N}.
\end{aligned}
\end{equation*}
Substituting this value into equation~\eqref{genI}, we obtain an exact solution of the model, given by
\begin{equation}
\begin{aligned}
I(t) &= \frac{N \left(\frac{I_0}{I_0 - N}\right) e^{N \beta t}}{\left(\frac{I_0}{I_0 - N}\right) e^{N \beta t} - 1}.
\label{ExactI1}
\end{aligned}
\end{equation}
With some algebraic manipulation, we can obtain the exact solution presented in class. We observe:
\begin{equation} \label{eq:solution}
\begin{aligned}
I(t) &= \frac{N I_0 e^{N \beta t}}{I_0 e^{N \beta t} - I_0 + N} \\
&= \frac{N I_0 e^{N \beta t}}{I_0 \left(\frac{N}{N}\right) e^{N \beta t} - I_0 \left(\frac{N}{N}\right) + N} \\
&= \frac{N I_0 e^{N \beta t}}{N\left(\left(\frac{I_0}{N}\right)(e^{N\beta t} - 1) + 1\right)} \\
I(t) &= \frac{I_0 e^{N \beta t}}{1 + \frac{I_0}{N}(e^{N \beta t} - 1)}.
\end{aligned}
\end{equation}

When $I_0 = 0$, we have $I(t) = 0$. Likewise, when $I_0 = N$, we have $I(t) = N$. Assume $I_0 \in (0, N)$.
Taking the limit of equation~\eqref{ExactI1} as $t \to \infty$, we observe the behaviour of any solution $I(t) \in (0,N)$ as $t$ gets arbitrarily large. Dividing both numerator and denominator in equation~\eqref{eq:solution} by $e^{N \beta t}$, we have
\begin{equation*}
\begin{aligned}
\lim_{t \to \infty} I(t) &= \lim_{t \to \infty}  \frac{I_0}{ e^{-N \beta t} + \frac{I_0}{N}(1 -  e^{-N \beta t})} = N.
\end{aligned}
\end{equation*}
As a result of the above limit, we can conclude that for all initial conditions $I(0) \in (0,N)$, the solution ultimately converges to the endemic equilibrium given by $I_\ast = N$.
Therefore, the endemic equilibrium is globally asymptotically stable.
      \end{proof}
    }
    
  \item \SIanalQbii
    
    {\color{blue}
      \begin{proof}
        Denote
\begin{equation*}
\frac{dI}{dt} = \beta I (N-I) = F(t, I).
\end{equation*}
Observe that $F(t, I)$ is $C^1$ for all $(t, I) \in \mathbb{R} \times [0, N]$. 
By the Fundamental Existence and Uniqueness Theorem, there exists a unique solution through any initial point $(0, I(0)) \in \{ 0\} \times [0, N]$.
In particular, since $I(t) = 0$ and $I(t) = N$ are the only equilibrium solutions, $I(0) \in (0, N)$ implies $I(t) \in (0, N)$ for all $t \in \mathbb{R}$.

Note that
\begin{equation*}
F(t, I) = \beta I (N-I) > 0
\end{equation*}
when $(t, I(t)) \in \mathbb{R} \times (0, N)$. 
So $I(t)$ is a monotonically increasing function of $t$.
To yield contradiction, given $\epsilon > 0$, suppose there does not exist $t < \infty$ such that $I(t) \in [N - \epsilon, N)$ for some $I(0) \in (0, N)$.
In other words, $I(t) \leq N-\epsilon$ for all $t \in \mathbb{R}$.
Since $I(t)$ is a monotonically increasing function and is bounded, it converges to its supremum, which is less than or equal to $N-\epsilon$. So
\begin{equation*}
\lim_{t \to \infty} I(t) = \sup_{t \in \mathbb{R}} I(t) = \hat{I} \implies \lim_{t \to \infty} F(t, I) = \lim_{I \to \hat{I}} F(t, I) = 0.
\end{equation*}
Since $\hat{I} \in (0, N)$,
\begin{equation*}
F(t, \hat{I}) > 0.
\end{equation*}
This contradicts our previous observation that $F(t, I)$ is a $C^1$ function. 

Given $\epsilon > 0$, there exists $t < \infty$ such that for any $I(0) \in (0, N)$, $I(t) \in [N - \epsilon, N)$. 
This implies that for any $I(0) \in (0, N)$,
\begin{equation*}
\lim_{t \to \infty} I(t) = N.
\end{equation*}
Therefore, the equilibrium point $I=N$ is globally asymptotically stable.
      \end{proof}
    }
    
  \end{enumerate}
\end{enumerate}

\section{Analysis of the basic SIR  model}

\basicSIRanalIntro
\begin{enumerate}[(a)]
\item \basicSIRanalQa

{\color{blue}
\begin{proof}[Solution]
Peak prevalence is defined as the maximum proportion of the population that is simultaneously infected. 
When $I(0) = 0$, $dI/dt = 0$ for all $t > 0$. So $I(t) = 0$ for all $t > 0$ and peak prevalence will be 0 as well.

Suppose $I(0) > 0$.
Recall that the derivative at a local maximum is 0. Let $t_p$ denote the time at which peak prevalence occurs. Let ${I}_p$ and ${S}_p$ be proportion of infected and susceptible individuals in the population at $t_p$. Then, ${I}_p$ represents peak prevalence and the following equation holds.
\begin{equation}
\left.\frac{dI}{dt}\right|_{t_p} = \R_0 S_p I_p - I_p = 0
\label{eq:peakp}
\end{equation}
Since $I_p > 0$ for all $t>0$, solving equation~\eqref{eq:peakp} yields the following:
\begin{equation}
S_p = \frac{1}{\R_0}.
\label{eq:peaks}
\end{equation}
Recall that the analytical solution of the SIR model for the phase portrait is given by
\begin{equation*}
I - I_0 = -(S - S_0) + \frac{1}{\R_0} \log (S/S_0)
\end{equation*}
At $t_p$, we have
\begin{equation*}
I_p - I_0 = -(S_p - S_0) + \frac{1}{\R_0} \log (S_p/S_0)
\end{equation*}
By substituting equation~\eqref{eq:peaks}, we obtain the following expression for peak prevalence:
\begin{equation*}
I_p = I_0 - \left(\frac{1}{\R_0} - S_0\right) + \frac{1}{\R_0} \log \left( \frac{1}{\R_0 S_0}\right).
\end{equation*}

Public health officials should be interested in peak prevalence for two reaons.
First, it measures how dangerous disease can be. 
Disease with higher peak prevalence will spread more easily and potentially be more dangerous.
Knowing peak prevalence early in an outbreak allows public health officials to estimate how much intervention is required and make interventions before peak prevalence is reached. Public health officials are interested in whether they have enough resources (such as hospital capacity) to treat the number of infected individuals.
Second, knowing peak prevalence allows future epidemics to be predicted.
Once we know how fast an epidemic grows, it is possible to estimate how fast the peak prevalence will be reached. 
Although exact shape of an epidemic is difficult to predict, knowing when peak prevalence occurs will allow public health officials to figure out an approximate shape of an epidemic.
\end{proof}
}

\item \basicSIRanalQb
  \begin{enumerate}[(i)]
  \item \basicSIRanalQbi

{\color{blue}
\begin{proof}[Solution]
Combining the expression for $dS/dt$ and $dR/dt$, we have that
\begin{equation*}
\frac{1}{S} \frac{dS}{dt} = - \R_0 \frac{dR}{dt}.
\end{equation*}
Note that $d \log S/dt = (dS/dt)/S$. We can integrate over $[0, \tau]$ to obtain
\begin{equation}\label{eq:S(R)}
\begin{aligned} 
\int_0^\tau \frac{d\log S}{dt} &= - \R_0 \int_0^\tau \frac{dR}{dt}.\\
\log S(\tau) - \log S(0) &= - \R_0(R(\tau) -R(0)) \\
S(\tau) &= S(0) e^{-\R_0(R(\tau)-R(0))}
\end{aligned}
\end{equation}

We recall that $S+I+R=1$, and thus we substitute $S=1-I-R$ into the above expression. This yields the following:
\begin{equation*}
1-I(\tau)-R(\tau)=S(0)e^{-\R_0(R(\tau)-R(0))}
\end{equation*}

If we substitute $I = dR/d\tau$, then 
\begin{equation*}
\begin{aligned}
\frac{dR}{d\tau}&=1-S(0)e^{-\R_0(R(\tau)-R(0))}-R(\tau) \\
\int_0^t d\tau&=\int_{R(0)}^{r} \frac{1}{1-S(0)e^{-\R_0(R-R(0))}-R} dR.
\end{aligned}
\end{equation*}
where $r$ is proportion of individuals recovered at time $t$, i.e. $r= R(t)$.

Thus, we have derived an expression for $t(r)$, given by Equation \eqref{eq:t(R)}. We note that it cannot be simplified beyond this form, as there isn't an exact solution for this integral.
\begin{equation}
\label{eq:t(R)}
t(r)=\int_{R(0)}^{r} \frac{1}{1-S(0)e^{-\R_0(R-R(0))}-R} dR
\end{equation}
\end{proof}
}

 \item \basicSIRanalQbii

{\color{blue}
\begin{proof}[Solution]
We start from the expression ${S(R)} = S_0 e^{-\R_0(R-R_0)}$ derived in equation~\eqref{eq:S(R)}. We can obtain an expression for $R(S)$ in the following way:
\begin{equation} \label{eq:R(S)}
\begin{aligned}
e^{-\R_0(R-R_0)} &= \frac{S}{S_0}\\
-\R_0(R-R_0) &= \log\left(\frac{S}{S_0}\right) \\
R(S) &= R_0 - \frac{1}{\R_0} \log\left(\frac{S}{S_0}\right)
\end{aligned}
\end{equation}

Recall from Equation (\ref{eq:peaks}) that at the time of peak prevalence $t_p$, the proportion susceptible in the population is
\begin{equation*}
\begin{aligned}
S_p := S(t_p) = \frac{1}{\R_0}.
\end{aligned}
\end{equation*}

Since $S(t)$ is monotonically decreasing and $R(t)$ is monotonically increasing (for non-equilibrium solutions), the values of $S(t)$ and $R(t)$ at $t_p$ are unique, and we can therefore substitute $S =S_p$ in (\ref{eq:R(S)}) to yield
\begin{equation*}
\begin{aligned}
R_p &= R_0 - \frac{1}{\R_0}\log\left(\frac{1}{S_0 \R_0}\right).
\end{aligned}
\end{equation*}

To obtain $t_p$, we can evaluate the expression for $t(R)$ given in (\ref{eq:t(R)}) at $R = R_p$ to obtain the time of peak prevalence. An expression for $t_p$ is therefore given by 
\begin{equation*}
\begin{aligned}
t_p = t(R_p)=\int_{R_0}^{R_p} \frac{1}{1-S_0e^{-\R_0(R-R(0))}-R} dR, 
\end{aligned}
\end{equation*}
where 
$$
R_p = R(0) - \frac{1}{\R_0}\log\left(\frac{1}{S_0 \R_0}\right).
$$

\end{proof}
}

  \item \basicSIRanalQbiii

{\color{blue}
\begin{proof}[Answers]
In order to compare our expressions with the P\&I time series for Philiadelphia 1918, we would have to make assumptions about how our model is related to the mortality data. 
In our basic SIR model, $R$ measures the proportion of the population ``removed'' from the dynamics of the system, either through death or recovery (and subsequent immunity). A relationship between $R(t)$ and the observed Philadelphia mortality curve could be established by introducing a new death rate parameter $0 \leq \delta \leq 1$ that determines the proportion of individuals that enter the removed class through death.
Then, $\delta R$ would represent the proportion of the population that have died from disase.
The associated assumption with $\delta$ would be that a constant proportion of infected individuals die per unit time. That is, every infected individual has an identical probability of death. Furthermore, we would be making the assumption that as soon as an individual enters the infected class, the probability of death is immediately $\delta$, when in reality there is likely a time delay associated with the probability of death. In other words, we wouldn't expect an individual who has been infected for 10 seconds to have the same probability of death as one who has been infected for two days.

Note that $\delta R$ is a cumulative function that represents proportion of population that have died whereas P\&I time series represents daily reports on death cases. 
In order to make a comparison, we would first have to make assumptions about the population size, $N$, so that $\delta N R$ represents number of individuals that have died.
Then, $\delta N R (t_{n+1}) - \delta N R (t_n)$ represents number of individuals that die between time $t_n$ and $t_{n+1}$. 
We can compare this quantity with daily reports by letting $t_n$ be in the unit of days.
To be more realistic, we can also assume that there is under reporting of death (only $\rho$ proportion of death is reported) and compare $\rho \delta N (R(t_{n+1}) - R(t_n))$ with time series.

Additionally, our basic SIR model also assumes a homogeneous population in which all susceptible individuals have the same probability of infection, and all infected individuals have the same probability of death. It is likely this was not the case with the 1918 influenza epidemic. This assumption could in theory be made weaker by incorporating age structure into the model.

We would not incorporate this expression in a report for a public health agency because of some key assumptions in the model that do not reflect the nature of an influenza epidemic. In particular, the temporal delay between prevalence and death is a factor we would expect to have a large impact on the inference of the time peak prevalence from mortality data, and the lack of heterogeneous age mortality structure may make our model an unreliable proxy for P\&I mortality data. This leads us to conclude that there would be large error associated with an estimate of the time of peak prevalence $t_p$ based on mortality data, and would not be useful for practical purposes.
\end{proof}
}

  \item \basicSIRanalQbiv

{\color{blue}
\begin{proof}[Answers]
No, it is not possible to find an exact analytic expression for $t(S)$. We recall that our expression for $t(R)$ could not be integrated using elementary functions - we needed to leave this expression as an integral. It follows then, that since ${S(R)} = S_0 e^{-\R_0(R-R_0)}$, $t(S) = t(S(R))$ also cannot be evaluated. This is because $S$ is equal to a function of $R$ which is written using elementary functions. Thus, $t(S)$ is simply a change of variables from $t(R)$. Therefore, we conclude that just as we could not evaluate the integral for $t(R)$, we similarly cannot evaluate the integral for $t(S)$ using elementary functions. Instead, we derive the below integral expression for $t(S)$ using $R(S)=R_0-\frac{1}{\R_0} \log \left( \frac{S}{S_0}\right)$ and $dS/dt=-\R_0SI$. 
\begin{equation}
\begin{aligned}
\frac{dS}{dt}&=-\R_0 SI\\
&=-\R_0 S(1-S-R)\\
&=-\R_0S \left( 1-S-R_0 + \frac{1}{\R_0} \log \left( \frac{S}{S_0}\right) \right)
\end{aligned}
\end{equation}
Thus, we get Equation (\ref{eq:t(S)}) as our final expression for t(S).
\begin{equation}
\label{eq:t(S)}
t(S)=-\frac{1}{\R_0}\int_{S(0)}^{S(\tau)} \frac{1}{1-S-R_0 +\frac{1}{\R_0}\log(\frac{S}{S_0})} dS
\end{equation}
\end{proof}
}

  \end{enumerate}
\item \basicSIRanalQc

{\color{blue}
\begin{proof}[Answers]
First, consider that for any initial condition $(S_0, I_0)$ with $S_0 > 1/\R_0$ and $I_0 >0$, we have $dS/dt < 0$, and $dI/dt > 0$. Since $dI/dt \geq 0$ as long as $S \geq 1/\R_0$, we have that $I >0$. It follows this, that as $dS/dt=-\R_0SI$, then $dS/dt < 0$ for $S \geq 1/\R_0$. That is, $S(t)$ is strictly monotonically decreasing in this region. Furthermore, it follows from the continuity of $dS/dt$ that for any $ \R_0 > 1$, $S(t) < 1/\R_0$ for some $t < \infty$.

%First, we will show that the set $A = \{(s, 0) \, | \, 0 \leq s \leq 1 \} \subseteq \Delta$ is a closed subset of $\Delta$. We will show prove this by show that $\Delta \setminus A$ is open. Consider any $(x, y) \in \Delta \setminus A$.  Then $y > 0$, and $d((x, y), A) = d((x, y), (0, y)) = y$, where $d$ is the usual (Euclidean) metric on $\mathbb{R}^2$. Define $U = B((x, y), y/2) \cap \Delta$, where $B((x, y), r)$ denotes the ball of radius $r$ about $(x, y)$. \\

%It follows then that the $U$ is an open set containing $(x, y)$ such that $U \cap A = \emptyset$, and $A$ is therefore closed.

Let $\Delta$ denote the biologically relevant region $ \{ (S, I) \, | \, S \geq 0, I \geq 0, 0 \leq S+I \leq 1 \}$, and $A$ denote the closed subset of $\Delta$ given by $[0, 1] \times \{0\}$. Note that based on the above result, to show $I =0$ is globally asymptotically stable, it is sufficient to consider initial conditions in $\Delta' = \{ (S, I) \, | \, 0 \leq S \leq 1/\R_0, I \geq 0, 0 \leq S+I \leq 1 \}$, since any non-equilibrium (i.e. $I_0 \neq 0$) trajectory in $\Delta$ will enter $\Delta'$ in finite time.

Next, we consider the function $L(S,I)=I$, and demonstrate that it is a strict Lyapunov function on $\Delta' \setminus A$. To begin, we check that $L$ is positive definite on $\Delta'$.
On $A$, we have that $I = 0$. It follows directly that $L(x)= I = 0 \ \forall \ x \in A$. Similarly, for any point $x$ in the biologically relevant region $\Delta' \setminus A$, $L(x) = I > 0$ (otherwise $x$ would be in $A$ by construction). Thus, $L(S,I)=I \geq 0 \ \forall \ (S,I) \in \Delta'$ and $L(S,I)=I > 0 \ \forall \ (S,I) \in \Delta' \setminus A$. This demonstrates that $L$ is positive definite on $\Delta'$.

Next, we show that $\dot{L}(x) < 0 \ \forall \ x \in \Delta' \setminus A$. First, let us find an expression for $\dot{L}(x)$. 
\begin{equation*}
\begin{aligned}
\dot{L}=\frac{d}{dt}L(S,I) &= \frac{dL}{dI} \frac{dI}{dt}+ \frac{dL}{dS}\frac{dS}{dt}\\
\dot{L}&= \R_0 SI - I
\end{aligned}
\end{equation*}
Thus, $\dot{L} < 0$ if $\R_0 S < 1$. We note that $\R_0 >1$ in biologically relevant contexts - as this means that the disease is capable of causing an epidemic. Thus, $\R_0 S<1$ when $S< \frac{1}{\R_0}$. Therefore, the function $L(S,I)=I$ is a strict Lyapunov function on $\Delta' \setminus A$ (i.e. when $S< \frac{1}{\R_0}$), and we conclude that $I = 0$ is globally asymptotically stable.
\end{proof}
}

\item \basicSIRanalQd

{\color{blue}
\begin{proof}[Solution]
We note that the equilibria of the SIR model are located where $dI/dt=dS/dt=dR/dt=0$. It is sufficient to analyze just where $dI/dt=dS/dt=0$. Thus, we must solve the following system of equations:
\begin{equation*}
\begin{aligned}
0&=-\R_0 SI\\
0&= \R_0 SI - I
\end{aligned}
\end{equation*}
We find a continuum of equilibria given by $(S,I)=(S_0,0)$ for $S_0 \in [0,1]$.

We begin our stability analysis by demonstrating that the function $L(S,I)=I$ is a Lyapunov function. We use the definition of $A$ as given above and define $\Delta'' = \{ (S, I) \, | \, 0 \leq S \leq 1/\R_0, I \geq 0, 0 \leq S+I \leq 1 \}$. Following the above argument exactly, we establish that $L(S,I)$ is positive definite on $\Delta''$. Next, we demonstrate that $\dot{L}(x) \leq 0 \ \forall \ x \in \Delta'' \setminus A$. As $\dot{L}=\R_0 S I -I$, we find that $\dot{L} \leq 0$ when $\R_0S \leq 1$. Thus, as we are considering only when $\R_0>1$, we find that $L(S,I)=I$ is a Lyapunov function on $\Delta'' \setminus A$ (i.e. when $S \leq 1/\R_0)$. Thus, by Lyapunov's direct method, the equilibrium points given by $(S_0, 0)$ for $S_0 \in [0, 1/\R_0]$ are stable. 

Next, we choose a point $(S_\ast,I_\ast)=(x_1 ,0)$, where $x_1=1/\R_0 + 2\varepsilon$ for $\varepsilon >0$. At time $t_0$, we perturb the point by $\delta >0$ such that $(S(t_0),I(t_0))=(x_1 ,\delta)$. However, we recall from part 2 (c), that $S(t)$ is monotonically decreasing (as $dS/dt<0$). Thus, it follows from the continuity of $dS/dt$ that $S(t_1) < 1/\R_0$ for some $t_1 < \infty$. Thus for some $t_1$, $(S(t_1),I(t_1))=(x_2 ,0)$ where $x_2 < 1/\R_0$. Therefore, $||(S(t_0), I(t_0))-(S_\ast, I_\ast)||<\delta$ implies $||(S(t_1), I(t_1))-(S_\ast, I_\ast)||>2\varepsilon >\varepsilon$.  As this holds for all $\delta>0$ arbitrarily small, we have $\exists \ \varepsilon>0$ such that $\nexists \ \delta(\varepsilon,t)>0$ such that $||(S(t_0), I(t_0))-(S_\ast, I_\ast)||<\delta$ implies $||(S(t), I(t))-(S_\ast, I_\ast)||<\varepsilon \ \forall \ t>t_0$. Thus, the equilibrium points given by $(S_0,0)$ for $S_0 \in \left(1/\R_0, 1\right]$ are unstable.
\end{proof}
}

\end{enumerate}

\newpage
\section*{Notes on Lyapunov functions}\hypertarget{NotesLyapFuns}{}

\NotesOnLyapunovFunctions

\bibliographystyle{vancouver}
\bibliography{4mba1_2018}

\bigskip

\centerline{\bf--- END OF ASSIGNMENT ---}

\bigskip
Compile time for this document:
\today\ @ \thistime

\end{document}
